\documentclass[10pt,a4paper]{amsart}
\usepackage[margin=19mm]{geometry}
\usepackage{amsmath,amssymb,mathtools}
\usepackage[scr=rsfs,cal=euler]{mathalfa}
\usepackage{xfrac}
\usepackage[unicode,hidelinks,bookmarks=false]{hyperref}
\newcommand{\ie}{i.e.\ }
\newcommand{\cf}{cf.\ }
\newcommand{\resp}{respectively\ }
\newcommand{\Subsection}{\subsection}
\providecommand{\nobreakdash}{\nobreak\hbox{-}}
\setlength{\emergencystretch}{2em}
\multlinegap0pt
\usepackage{bm}
\usepackage{bbm}
\usepackage{dsfont}
\usepackage{xspace}
\usepackage{cancel}
\usepackage{mathtools}
\usepackage{amsthm}
\makeatletter
\@ifundefined{lem}{\newtheorem{lem}{Lemma}}{}
\@ifundefined{maintheorem}{\newtheorem{maintheorem}{Theorem}}{}
\makeatother
\title[On the dynamics of Halley's method]{On the dynamics of Halley's method}
\author{Gang Liu, Soumen Pal,, Saminathan Ponnusamy}
\date{Source: arXiv:2508.00631v2 (CC BY 4.0)}
\begin{document}
\maketitle

\noindent\emph{Source and licence.} On the dynamics of Halley's method. Version arXiv:2508.00631v2 (CC BY 4.0), distributed under the Creative Commons Attribution 4.0 International licence, as recorded in the arXiv licence metadata for this exact version and shown on the abstract page https://arxiv.org/abs/2508.00631v2. Full rights notice: licenses/arxiv-2508.00631-CC-BY-4.0.txt.\par\smallskip

\noindent\textit{Editorial scope.} Counted unit: the Maintheorem whose source label is \texttt{exact2} in the article (source0.tex lines 229--231, source label \texttt{exact2}), delivered together with the article's own complete original proof of it (source0.tex lines 232--259, one \texttt{proof} environment of 28 lines). No mathematics is written, repaired or re-derived: statement and proof are the article's own text line for line. Required context retained uncounted: bdry\_pole (lines 201--203); all\_pole (lines 205--207); gen\_dyn (lines 209--211). Every definition, display and label of those lines is retained unchanged. Omitted from this package and not redistributed: 1--200, 204--204, 208--208, 212--228, 260--746. Nothing inside a retained excerpt is omitted or altered: every retained body is delivered exactly as the article's lines are written, so no omission receipt of any kind accompanies this package. Citations: the retained excerpts carry no citation command at all, so no bibliography entry of the article is transcribed and no reference is redistributed. Reference adaptation: none was needed and none was made; every reference inside the delivered unit resolves inside it, so no unresolved reference or citation is produced. Printed slips kept literal: no printed word, symbol, hypothesis or number of the counted statement or of its proof was altered. Layout and class: the article's own class is \texttt{article} with packages including bbm, graphicx, wrapfig, euscript, subcaption, hyperref, epstopdf, enumerate, amsmath, amsfonts, amssymb, amsthm, epsfig, titling; the wrapper is the lane's pinned \texttt{amsart} editorial wrapper at 10pt on A4, which restates the theorem-like environments the retained text needs and adds the article's own macro definitions that the retained text needs (none), with extra wrapper packages (bm, bbm, dsfont, xspace, cancel, mathtools). The delivered numbering is the unit's own. Assets: no retained span contains a figure environment, a graphics-inclusion command, a PDF-page inclusion, a table or a TikZ picture; no image file is copied into this package, the package declares no asset dependency and no image file is redistributed with it. Licence: version arXiv:2508.00631v2 of this article is distributed under the Creative Commons Attribution 4.0 International licence, as recorded in the arXiv licence metadata for this exact version and shown on the abstract page https://arxiv.org/abs/2508.00631v2. A full rights notice is delivered at licenses/arxiv-2508.00631-CC-BY-4.0.txt. Characters: every retained excerpt is pure ASCII, so the delivered engine, XeLaTeX with its default Latin Modern text font, reports no missing character.
\begin{lem}\label{bdry_pole}
The boundary of an unbounded immediate attracting basin of $H_p$ corresponding to a root of $p$ contains at least one pole.
\end{lem}

\begin{lem}\label{all_pole}
If there is an unbounded Julia component of $H_p$ containing all the poles, then $\mathcal{J}(H_p)$ is connected.
\end{lem}

\begin{lem}\label{gen_dyn}
Suppose $R$ is a rational function with real coefficients having two real fixed points $x_{1}$ and $x_{2}\left(x_{1}<x_{2}\right)$ such that $R$ does not have any critical point, pole, or fixed point on $(x_{1}, x_{2})$. Then  the successive iterations $\{R^n(x)\}$ of any point $x$ in $(x_{1}, x_{2})$ will converge to either $x_{1}$ or $x_{2}$. More precisely, if $R(x)<x$ (or $R(x)>x$) then $\lim\limits_{n \to \infty} R^n(x)=x_1$ (or $\lim\limits_{n \to \infty} R^n(x)=x_2$, respectively).
\end{lem}

\begin{maintheorem}\label{exact2}
	The Julia set of $H_p$ is a line if and only if $p$ has exactly two distinct roots, each with the same multiplicity.
\end{maintheorem}

\begin{proof}
First we consider that $p$ is having exactly two roots with same multiplicity. Using the Scaling property we can consider $p(z)=\left(z^{2}-1\right)^{k}$. Then	$$
H_{p}(z)=\frac{z\left((2 k-1) z^{2}+3\right)}{(2 k+1) z^{2}+1}\qquad\text{and}\qquad H_{p}^{\prime}(z)=\frac{\left(4 k^{2}-1\right) z^{4}-6 z^{2}+3}{\left((2 k+1) z^{2}+1\right)^{2}}.
$$
The free critical points and the poles of $H_p$ are the solutions of
$ z^{2}=\frac{3 \pm 2i \sqrt{3} \sqrt{k^{2}-1}}{\left(4 k^{2}-1\right)}$ and $z^2=- \frac{1}{2k+1}$ respectively, which are non-real complex numbers.
As $\mathcal{J}(H_{p})$ is symmetric about both the axes, all critical points are in $\cup \mathcal{A}_{ \pm 1}$. As the origin is the only finite extraneous fixed point, it follows from Lemma~\ref{gen_dyn} that
$\mathcal{A}_{ \pm 1}$ are unbounded, and hence, their respective boundary contains at least one pole (by Lemma~\ref{bdry_pole}). Again, as poles are non-real, $\partial \mathcal{A}_{ \pm 1}$ contains both the poles. Hence, by Lemma~\ref{all_pole} we get that
$\mathcal{A}_{\pm 1}$ are simply connected. Using Riemann-Hurwitz formula, we get that the local degree of $H_p$ in each immediate basin is $3$, and therefore both the immediate basins are completely invariant. Hence, the Fatou set $\mathcal{F}(H_p)$ is the union of two immediate basins, and consequently, $\mathcal{J}(H_p)$ is a Jordan curve. Since the imaginary axis is invariant under the iteration of $H_{p}$ and both the poles are purely imaginary, as well as $\mathcal{J}(H_p)$ is preserved by the reflection about both the axes, we conclude that the Julia set is the imaginary axis.
\par
Conversely, let $\mathcal{J}(H_p)$ be a line.
Then $ p$ has exactly two roots, say $a$ and $b$, with multiplicity $k$ and $m$, respectively. Without loss of generality we can consider that $a$ and $ b$ are real.
Then $\mathcal{J}(H_p)$ is symmetric about the reflexion of real axis, and hence, $\mathcal{J}(H_{p})$ is a vertical line.
Again, using the Scaling property, we can consider the $\mathcal{J}(H_{p})$ as the imaginary axis. Note that the extraneous fixed point of $H_p$ is $\frac{b k+a m}{k+m}$, which is real and repelling. As $0 \in \mathcal{J}(H_p)$ is the only real number on the Julia set of $H_p$, we have
$
\frac{b k+a m}{k+m}=0$. This gives $b=-\frac{a m}{k} .
$
Therefore, we get that $ p(z)=(z-a)^{k}\left(z+\frac{a m}{k}\right)^{m} .
$ Consider the affine map $T(z)=a z$. Then $\mathcal{J}\left(H_{p}\right)=T\left(\mathcal{J}\left(H_{p}\right)\right)$, and hence, using the Scaling property we can consider
$
p(z)=(z-1)^{k}\left(z+\frac{m}{k}\right)^{m} .
$
Then $$H_{p}(z)=\frac{z\left(k(k+m-1) z^{2}+2(k-m) z+3 m\right)}{k(k+m+1) z^{2}+m}.$$
Note that, as $\mathcal{J}(H_{p})$ is the imaginary axis, therefore for every $y \in \mathbb{R}$, there exists $y^{*} \in \mathbb{R}$ such that
$H_{p}(i y)=i y^{*}.$
This can only be true if $k=m$.
Hence the result.
\end{proof}

\end{document}
